\begin{textsample}{POS Dim 9 – global\_south\_2025\_02 – Score 2.00 – global\_south\_2025\_02/120944401.txt}  \label{ex:f9_pos_256}
US President Donald Trump said on Thursday that Israel would hand over Gaza to the United States after fighting was over and the enclave ' s population was already resettled elsewhere, which he said meant no US troops would be needed on the ground.
A day after worldwide condemnation of Trump ' s announcement that he aimed to take over and develop the Gaza Strip into the"Riviera of the Middle East", Israel ordered its army to prepare to allow the"voluntary departure"of Gaza ' s residents.
Trump, who had previously declined to rule out deploying US troops to Gaza, clarified his plans in comments on his Truth Social web platform."The Gaza Strip would be turned over to the United States by Israel at the conclusion of fighting,"he said.
Palestinians"would have already been resettled in far safer and more beautiful communities, with new and modern homes, in the region.""No soldiers by the US would be needed!"he said. @ @ @ @ @ @ @ @ @ @ what Prime Minister Benjamin Netanyahu called Trump ' s"remarkable"proposal, Defence Minister Israel Katz said he had ordered the army to prepare a plan to allow residents who wished to leave to exit Gaza voluntarily."I welcome President Trump ' s bold plan, Gaza residents should be allowed the freedom to leave and emigrate, as is the norm around the world,"Katz said on X.
Katz said his plan would include exit options via land crossings, as well as special arrangements for departure by sea and air.
Trump ' s unexpected announcement on Wednesday, which sparked anger around the Middle East, came as Israel and Hamas were expected to begin talks in Doha on the second stage of a ceasefire deal for Gaza, intended to open the way for a full withdrawal of Israeli forces and an end to the war.
Regional heavyweight \textbf{Saudi} Arabia rejected the proposal outright and Jordan ' s King Abdullah, who will meet Trump at the White House next week, said on Wednesday he rejected any attempts @ @ @ @ @ @ @ @ @ @ will not sell our land for you, real estate developer.
We are hungry, homeless, and desperate but we are not collaborators,"said Abdel Ghani, a father of four living with his family in the ruins of their Gaza City home."If ( Trump ) wants to help, let him come and rebuild for us here."Displacement of Palestinians has been a long-standing and sensitive issue in the region, and many view Trump ' s proposal as an attempt to continue policies of forced relocation, which has been banned under the 1949 Geneva Conventions.
Some Israeli politicians, including former general Giora Eiland, supported the idea of relocation, viewing it as"logical,"but this view has been heavily criticized by others.
Israel ' s military campaign in Gaza has resulted in tens of thousands of deaths and forced many Palestinians to seek safety within the enclave.
However, there is widespread fear of permanent displacement, recalling the"Nakba"of 1948 when hundreds of thousands of @ @ @ @ @ @ @ @ @ @ Israel.
In response to international criticism, Katz suggested that countries such as Spain, Ireland, and Norway, which have opposed Israel ' s actions in Gaza, should take in displaced Palestinians as he argued that these countries, having made accusations against Israel, were"legally obligated"to offer refuge to Gaza ' s residents. ' I thought we voted for America first ' Trump ' s proposal for the United States to take control of the war-torn Gaza Strip has sparked confusion and skepticism among some Republicans, while others have voiced support for his"bold, decisive"approach.
The suggestion has drawn widespread international condemnation, along with dissent from certain Republican lawmakers, who have largely backed Trump ' s previous initiatives, including pausing foreign aid and reducing federal workers. \textbf{Skeptical} lawmakers emphasized their support for a long-standing two-state solution for Israel and the Palestinians, which has been a cornerstone of US diplomacy.
Some voiced opposition to using US taxpayer funds or deploying troops to a region devastated by over a @ @ @ @ @ @ @ @ @ @ for America first,"remarked Republican Senator Rand Paul on X.com."We have no business contemplating yet another occupation that will cost our treasure and blood."Republicans, who hold narrow majorities in Congress, have faced opposition from Democrats, with Senator Chris Van Hollen denouncing the proposal as"ethnic cleansing by another name."Republican Senator Jerry Moran echoed concerns, stating that the two-state solution can not be discarded unilaterally."It ' s not something that can be unilaterally decided,"he remarked.
Senator Lisa Murkowski expressed reluctance to speculate on the proposal, highlighting the risks of sending US forces to a region long beset by turmoil."I think that is quite frightening,"she said.
However, House Speaker Mike Johnson praised Trump ' s plan as a"bold, decisive action to try to secure the peace of that region.

% matched lemmas: saudi, skeptical
\end{textsample}
